% !PRINTED 208
\weberpage{208}
%<<CONT
Es ist hierbei immer angenommen, dass die Co\"efficienten
von $F(x_1, x_2)$ nicht alle gleich Null sind. Wir k\"onnen also,
wenn wir von dieser Voraussetzung absehen, den Satz 3. auch
so aussprechen:

4.\ \gesperrt{Ist $F(x_1, x_2)$ eine invariante Form der Gruppe $G$
von niedrigerem Grade als $G$, die sich nicht als
Product aus den Grundformen darstellen l\"asst,
so muss $F(x_1, x_2)$ identisch verschwinden.}

\rule{4cm}{0.4pt}
